% Цикл пересчёта для плана, построенного методом северо-западного угла.
% Старт — пустая клетка (1,4); знаки чередуются вдоль ломаной.
\begin{figure}[htbp]
\centering
\begin{tikzpicture}[scale=1.0, >={Stealth[length=2.2mm]}]

\def\cw{1.6}
\def\ch{1.1}

% сетка 3 x 4
\foreach \i in {0,1,2,3}
  \draw[blue!55!black] (0, -\i*\ch) -- (4*\cw, -\i*\ch);
\foreach \j in {0,1,2,3,4}
  \draw[blue!55!black] (\j*\cw, 0) -- (\j*\cw, -3*\ch);

% заголовки
\foreach \j/\name in {0/{B_1}, 1/{B_2}, 2/{B_3}, 3/{B_4}}
  \node[font=\footnotesize] at ({(\j+0.5)*\cw}, 0.3) {$\name$};
\foreach \i/\name in {0/{A_1}, 1/{A_2}, 2/{A_3}}
  \node[font=\footnotesize, anchor=east] at (-0.18, {-(\i+0.5)*\ch}) {$\name$};

% содержимое клеток: значения плана и крестики
\newcommand{\cell}[3]{\node[font=\small] at ({(#2+0.5)*\cw}, {-(#1+0.55)*\ch}) {#3};}
\cell{0}{0}{$27$} \cell{0}{1}{$\times$} \cell{0}{2}{$\times$} \cell{0}{3}{$\times$}
\cell{1}{0}{$6$}  \cell{1}{1}{$13$}     \cell{1}{2}{$1$}      \cell{1}{3}{$\times$}
\cell{2}{0}{$\times$} \cell{2}{1}{$\times$} \cell{2}{2}{$26$} \cell{2}{3}{$17$}

% ломаная цикла: (1,4) -> (1,1) -> (2,1) -> (2,3) -> (3,3) -> (3,4) -> (1,4)
\coordinate (c14) at ({3.5*\cw}, {-0.3*\ch});
\coordinate (c11) at ({0.5*\cw}, {-0.3*\ch});
\coordinate (c21) at ({0.5*\cw}, {-1.3*\ch});
\coordinate (c23) at ({2.5*\cw}, {-1.3*\ch});
\coordinate (c33) at ({2.5*\cw}, {-2.3*\ch});
\coordinate (c34) at ({3.5*\cw}, {-2.3*\ch});

\draw[orange!75!black, thick, ->] (c14) -- (c11);
\draw[orange!75!black, thick, ->] (c11) -- (c21);
\draw[orange!75!black, thick, ->] (c21) -- (c23);
\draw[orange!75!black, thick, ->] (c23) -- (c33);
\draw[orange!75!black, thick, ->] (c33) -- (c34);
\draw[orange!75!black, thick, ->] (c34) -- (c14);

% знаки в вершинах
\foreach \p/\s in {c14/{\oplus}, c11/{\ominus}, c21/{\oplus},
                   c23/{\ominus}, c33/{\oplus}, c34/{\ominus}}
  \node[circle, fill=white, draw=orange!75!black, inner sep=0.8pt,
        font=\footnotesize] at (\p) {$\s$};

\node[font=\footnotesize, anchor=north] at ({2*\cw}, {-3*\ch-0.3})
  {$\theta = \min\limits_{\ominus} x_{ij} = \min\{27,\ 1,\ 17\} = 1$};

\end{tikzpicture}
\caption{Цикл пересчёта для плана метода северо-западного угла. Стартовая клетка
$(1,4)$ пуста и получает знак $\oplus$; дальше знаки чередуются, повороты ломаной
допускаются только в заполненных клетках. Шаг $\theta$ равен наименьшему объёму
среди клеток со знаком $\ominus$; после пересчёта клетка $(2,3)$ обнуляется и
выходит из базиса.}
\label{fig:cycle}
\end{figure}
